\phantomsection
\chapter*{Abstract and indexation in English}\label{chap:abstract}
\addcontentsline{toc}{chapter}{Abstract and indexation in English}

Hybrid gender agreement in Polish

Abstract: 

The aim of this thesis is to research hybrid agreement, as described by Greville Corbett in his work \citep{corbett_agreement_2023, corbett_agreement_1979}, with a focus on gender in Polish. We conducted a corpus study using three types of nouns in Polish that can exhibit hybrid gender agreement: socially gendered nouns, profession nouns and depreciative forms. Based on 1722 agreement occurrences with those nouns found in the National Corpus of Polish we found that the majority of occurrences had syntactic agreement (74\%). We also found effects of target-controller ordering on the choice between semantic and syntactic agreement in the whole dataset. We also found an effect of distance in agreement occurrences with socially gendered nouns and an effect of type of target in occurrences with profession nouns. We then conducted two acceptability judgement experiments, one with profession nouns and the other with depreciative forms. We did not find significant effects of semantic agreement, but unlike in the corpus study, we found a preference for semantic agreement across almost all of the conditions. The studies show that semantic agreement is acceptable in Polish and sensitive to several factors of Corbett's Agreement Hierarchy. The results also give directions for further studies on the topic of hybrid gender agreement in Polish.
\\


English keywords: agreement; gender; Polish; corpus; experiments
\\


Publication type: thesis
\\

MeSH : Academic Dissertation 




